\subsection{平坦 $G$-丛模空间的量子化}\label{flatg}

设 $G$ 为紧 Lie 群，复化为 $G_\C$。记 $\g$ 与 $\g_\C$ 为相应的 Lie 代数。$G_\C$ 有反全纯对合 $\varrho$，其不动点集为 $G$。设 $C$ 为紧闭可定向二维流形，$M$ 为指定拓扑型的平坦 $G$-丛\footnote{本讨论中我们忽略 $M$ 的奇点带来的技术问题。例如，取 $G=\SO(3)$ 并要求丛 $E$ 满足 $w_2(E)\ne0$，即可完全避开这类问题。} $E\to C$ 的模空间。$M$ 有自然的辛结构 $\omega_0$，我们将其归一化，使最小非零周期为 $2\pi$。对 $G=\SU(N)$，它可由公式 \cite{ABott}
\be\label{defform}\omega_0=\frac{1}{4\pi}\int_C\Tr\,\delta A\wedge\delta A\ee
定义，其中 $\Tr$ 是在 $G$ 的基本 $N$ 维表示中取迹，$\delta A$ 表示 $E$ 上平坦联络的变分。周期条件意味着 $\omega_0$ 是线丛 $\L\to M$ 的曲率（该线丛生成 $M$ 的 Picard 群模挠）。人们希望用辛结构\footnote{稍后提到的 Chern--Simons 与 WZW 模型的“等级”，按通常定义并不是 $\kappa$。为简单起见设 $G$ 单连通，它是 $k=\kappa-h$，其中 $h$ 是 $G$ 的对偶 Coxeter 数。对单连通的 $G$，有 $K_M^{1/2}\sim\L^{-h}$。在某些处理中，自然地把 $\frL\otimes K_M^{1/2}$ 而非 $\frL$ 称为预量子线丛，公式也直接用 $k$ 而非 $\kappa$ 写出。}
$\omega=\kappa\omega_0$ 与预量子线丛 $\frL=\L^\kappa$ 来量子化 $M$，其中 $\kappa$ 为某个整数。假设 $\kappa>0$ 不失一般性，因为反转 $C$ 的定向会使 $\kappa$ 变号。

$M$ 是 $\R\times C$ 上规范群为 $G$ 的三维 Chern--Simons 规范理论的相空间，其中 $\R$ 参数化“时间”\cite{WittenJ}。量子化的目标是构造与空间流形 $C$ 相伴的 Hilbert 空间。$M$ 的辛结构只用到 $C$ 的定向，不用 $C$ 的其他结构。因此形式上可以希望，量子化能够以这样的方式进行，使 $M$ 的保定向微分同胚自然地作用在量子 Hilbert 空间上。事实上，若要从三维 Chern--Simons 理论得到拓扑场论，这一点是必需的。

对膜量子化而言，第一步是选取 $M$ 的一个好的复化。确实，$M$ 有好的复化 $Y$，即平坦 $G_\C$-丛 $\h E\to C$ 的模空间。它是复辛流形；复结构来自 $G_\C$ 的复结构，而全纯辛结构由公式 \eqref{defform} 把 $\g$-值平坦联络 $A$ 换成 $\g_\C$-值平坦联络 $\cA$ 来定义，是 $M$ 的辛结构的解析延拓。$Y$ 上典型的全纯函数，是 $\cA$ 沿 $C$ 中给定定向环的和乐在某个全纯表示中的迹。为与本文记号一致，我们称此复结构为 $I$；不过在 Hitchin 关于 Higgs 丛的工作 \cite{Hitchin}（其相关性稍后说明）以及该工作对几何 Langlands 的应用中，它通常被称为 $J$。$Y$ 有反全纯对合 $\tau$ 固定 $M$。$\tau$ 可通过把 $\varrho$ 作用在 $\cA$ 或其单值化上来定义。

$Y$ 的选择并未破坏 $C$ 的微分同胚对称性；$C$ 的保定向微分同胚作用在 $Y$ 上，并保持其复辛结构。因此形式上，膜量子化的程序给出 $M$ 量子化的微分同胚不变答案：它就是 $\H=\Hom_{A_K}(\B,\B_\cc)$。

我们可以通过选取全纯极化或选取超 Kähler 度量，试着把这一点具体化。（后一途径在 \cite{GW} 中有描述，其中某些细节比我们这里更详细。）两种途径都涉及一个看似破坏微分同胚对称性的辅助选择。但与 $\Hom_{A_K}(\B,\B_\cc)$ 的等价性表明，来自全纯极化或超 Kähler 度量的输出确实具有微分同胚对称性。

使 $M$ 的膜量子化更具体的两种途径，起点都是选取 $C$ 上的复结构。当然没有自然的选择；在连通到恒等的微分同胚之下，可能的选择构成 Teichmüller 空间 $\T$。对 $t\in\T$，记 $\I_t$ 为 $C$ 上相应的复结构。

给定 $\I_t$ 之后，就可以用 $Y$ 的全纯极化或超 Kähler 结构使膜量子化具体化。先描述全纯极化的途径。

给定 $\I_t$ 之后，平坦丛具有全纯结构（由平坦联络的 $(0,1)$ 部分定义）与反全纯结构（由 $(1,0)$ 部分定义）。为定义 $Y$ 的全纯极化，在 $Y$ 上放置一个等价关系：两个平坦 $G_\C$-丛若具有相同的全纯结构，就视为等价。$Y$ 中每点属于唯一的等价类。考虑包含平坦联络 $\cA=\d\bar z\cA_{\bar z}+\d z\cA_z$ 的等价类。同类中任何其他联络（模规范变换）有相同的 $(0,1)$ 部分，故形如
$\cA'=\cA+\W$，其中 $\W=\d z\W_z$ 是取值于 $\ad(\h E)$（$\h E$ 的伴随丛）的 $(1,0)$-形式。若 $\cA$ 平坦，则 $\cA'$ 也平坦的条件是 $\W$ 全纯，也就是说它取值于 $H^0(C,K_C\otimes\ad(\h E))$。指标定理表明该向量空间的维数为 $n=\frac12\dim_\C Y$。于是这些等价类是 $\C^n$ 的拷贝，此构造定义了 $Y$ 的全纯极化 $\varPi_t$。

可以这样看出 $\varPi_t$ 在 $M$ 上诱导出复极化。由 $\varPi_t$ 的定义，$M$ 的每点都含于唯一的叶；剩下的问题是证明每个叶至多含 $M$ 中的一个点。这一陈述是 Narasimhan--Seshadri 定理的一部分，该定理建立了 $C$ 上稳定全纯 $G_\C$-丛与平坦 $G$-丛之间的对应。交 $F\cap M$ 是给定全纯 $G_\C$-丛上平坦酉 $\g$-值联络的空间。Narasimhan--Seshadri 定理说，若 $E$ 稳定，则 $F\cap M$ 是单点；否则为空。

用复平坦联络的展开 $\cA=\cA_{\bar z}\d\bar z+\cA_z\d z$ 来看，叶空间由 $\cA_{\bar z}$ 全纯地参数化。限制到 $M$ 后，展开变为 $A=\d\bar zA_{\bar z}+\d zA_z$，其中 $A_z=-A_{\bar z}^\dagger$（这里假设约定为：酉平坦联络是反 Hermite 的）。因此 $\varPi$ 在 $M$ 上诱导的复结构，就是使 $A_{\bar z}$ 全纯变化的那一个。在此复结构（记作 $\J_t$）下，$M$ 可解释为 $C$ 上全纯 $G_\C$-丛的模空间。若写 $\delta A=\d z\delta A_z+\d\bar z\delta A_{\bar z}$，则公式 \eqref{defform} 关于 $\delta A_z$ 与 $\delta A_{\bar z}$ 是双线性的，这表明辛形式关于 $\J_t$ 为 $(1,1)$ 型，故 $\J_t$ 是 $M$ 的复极化。这与第 \ref{scbq} 节的一般论证相符。稍加小心便可证明，$\J_t$ 其实是 $M$ 的 Kähler 极化。

第 \ref{branecp} 节的论证适用于此例，它们表明膜量子化的 Hilbert 空间 $\H$ 作为向量空间就是
$\H=H^0_{\J_t}(M,K_M^{1/2}\otimes\frL)$。这一等价也附带一个不太显式的 $\H$ 的 Hilbert 空间结构描述，其中涉及一个全纯三角形。

把反全纯对合 $\tau$ 作用于 $\varPi_t$，得到另一个全纯极化 $\varPi'_t$。此极化的叶由具有相同反全纯结构的平坦 $G_\C$-丛组成。$\varPi'_t$ 在 $M$ 上诱导相反的复结构 $-\J_t$。一般地，如第 \ref{branecp} 节所解释，用复极化的几何量子化在复结构变号下不变。因此用 $\varPi'_t$ 的几何量子化给出与 $\varPi_t$ 等价的结果。我们将不再讨论 $\varPi'_t$。

使 $Y$ 的膜量子化具体化的另一条途径如下。一旦给二维流形 $C$ 指定复结构 $\I_t$，就可以表述并求解 Hitchin 方程 \cite{Hitchin}，使平坦 $G_\C$-丛 $\h E$ 具有 Higgs 丛的结构。空间 $Y$ 获得新的复结构 $J_t$，在其中它参数化 Higgs 丛；$Y$ 的复辛结构则延拓为超 Kähler 结构。（$J_t$ 在 Hitchin 的工作及后续大量文献中称为 $I$。）具体地说，Hitchin 方程涉及把 $\h E$ 的结构群从 $G_\C$ 约化到 $G$。相对于这一分解，复平坦联络 $\cA$ 分解为
$\cA=A+\i\phi$，其中 $A$ 是 $G_\C$-值平坦联络，$A=\d\bar zA_{\bar z}+\d zA_z$ 是 $G$-值联络，$\phi=\d\bar z\phi_{\bar z}+\d z\phi_z$ 是 $C$ 上伴随值的一形式。$A$ 与 $\phi$ 满足 Hitchin 方程。复结构 $J_t$ 就是使 $A_{\bar z}$ 与 $\phi_z$ 为全纯变量的那一个。

第 \ref{hyperm} 节的讨论表明，$\H$ 作为向量空间可等同于 $H^0_{J_t}(M,K_M^{1/2}\otimes\frL)$。但这一等同并不能给出 $\H$ 的 Hilbert 空间结构的描述。

解 Hitchin 方程得到的复结构 $J_t$ 定义在整个 $Y$ 上，而通过忘记平坦丛的反全纯结构得到的复结构 $\J_t$ 只定义在 $M$ 上。然而 $J_t$ 限制到 $M$ 时与 $\J_t$ 相同。为看出这一点，只需注意 $M$ 作为轨迹 $\phi=0$ 嵌入 $Y$。因此，使 $A_{\bar z}$ 与 $\phi_z$ 全纯的复结构 $J_t$，限制到 $M$ 上就是使 $A_{\bar z}$ 全纯的复结构 $\J_t$。

因此，由全纯极化得到的 $\H$ 描述，与解 Hitchin 方程得到的描述相同。

无论是全纯极化的描述，还是解 Hitchin 方程的描述，都没有说清 $H^0_{\J_t}(M,K_M^{1/2}\otimes\frL)$ 为何、或在何种意义下与 $t\in\T$ 的选择无关。然而这正是与膜量子化的关系所预言的。它也是三维 Chern--Simons 规范理论的拓扑不变性 \cite{WittenJ}，以及获得二维共形场论 WZW 模型的全纯分解的满意图像 \cite{KZ} 所需要的。通常的做法是把 $H^0_{\J_t}(M,K_M^{1/2}\otimes\frL)$ 视为 $\T$ 上某个向量丛的纤维，并在该向量丛上构造（投影）平坦联络。积分这个平坦联络，就给出不同 $t$ 的 $H^0_{\J_t}(M,K_M^{1/2}\otimes\frL)$ 之间的等价，从而恢复拓扑不变性，并在二维共形场论与三维拓扑场论中得到满意的结果。该平坦联络已从多个观点构造出来 \cite{ADW,Hitchtwo,Faltings}；不同途径被证明给出等价的结果 \cite{And}。这个平坦联络与 WZW 共形场论的共形块空间上的自然平坦联络一致 \cite{Las}。为与三维 Chern--Simons 规范理论的酉性以及二维 WZW 模型的预期性质相容，这个平坦联络必须是酉的；也就是说，它必须保持纤维上某个正 Hermite 内积。于是平坦联络的整体单值化给出映射类群的酉表示（更确切地说，是其中心扩张，因为该联络只是投影平坦的）。已知的规范论构造并没有给出像几何量子化所建议的公式 \eqref{ngg} 那样简单的东西。人们相信 \cite{El,Gaw,Gaw2}，正确的公式需要一个额外因子，其形式为
\be\label{nggo}(\psi,\chi)=\int_M\d\mu\,\bar\psi\chi\,e^{G_\kappa},\ee
其中 $e^{G_\kappa}$ 是靶为 $G_\C/G$、等级为 $\kappa$ 的 WZW 模型的配分函数。这个因子来自从所有联络的空间约化到平坦酉联络空间 $M$ 时，对 $G_\C$-值规范变换轨道的积分。抽象地，已经证明该平坦联络实际上是酉的 \cite{And2}；更具体地，已知酉结构可由诸如 \eqref{nggo} 的公式定义，其中当 $\kappa$ 很大时有渐近展开 $G_\kappa=F+\sum_{n=1}^\infty\kappa^{-n}F_n$，首项 $F$ 是 $M$ 的 Ricci 势 \cite{And3}。在首阶，这一结果与基于 $G_\C/G$ 的 WZW 模型的预言一致 \cite{And4}。

膜量子化给出的 Hilbert 空间内积答案，就是我们在第 \ref{branecp} 节描述的那个。这一描述不太显式，因为它依赖于一个难以计算的全纯三角形。理解预期因子 $e^{G_\kappa}$ 如何在这一途径中出现，将会极其有趣。
